\begin{textsample}{POS Dim 1 – vem\_ed\_25 – Score 14.00 – vem\_ed\_25/t135.txt}  \label{ex:f1_pos_072}
5º Congresso LinFE: experiências dos PCIs A equipe dos PCIs / Cesu e duas professoras \textbf{que} realizam PCIs na Fatec Guaratinguetá – Patrícia Januária Barbosa e Tálita Guarino – apresentaram trabalhos no 5º Congresso de Línguas para Fins Específicos (LinFE).
O evento, focado na área de Linguística Aplicada voltada à formação, capacitação e atualização de professores de Línguas para Fins Específicos, ocorreu na modalidade híbrida nos dias 12 e 13 de setembro de 2024 e \textbf{foi} organizado pela Universidade Estadual do Rio de Janeiro (UERJ).
Patrícia Januária Barbosa e Tálita Guarino \textbf{compartilharam} suas experiências no trabalho intitulado “Uso das tecnologias e ensino de inglês para fins \textbf{específicos}: uma experiência de internacionalização (COIL) entre Fatec Guaratinguetá e Yerevan University”.
Ana Carolina Freschi, Priscilla Ferro e Regiane Camargo Moreira (equipe dos PCIs / Cesu) fizeram a apresentação “Intercâmbios Virtuais 101: a experiência dos Projetos Colaborativos Internacionais no Centro Paula Souza”.
As autoras falaram sobre as características e as fases do PCI e apontaram as \textbf{competências} desenvolvidas nesses projetos de Intercâmbio Virtual do tipo COIL.
Regiane Camargo Moreira fez a palestra “A abordagem LinFE na experiência dos Intercâmbios Virtuais”.
Mostrou o caminho do \textbf{desenvolvimento} dos projetos de Intercâmbio Virtual, \textbf{que} envolvem o planejamento de temáticas, a exploração de ferramentas para a colaboração on-line, os \textbf{encontros} virtuais entre docentes (desde o planejamento até a avaliação final) e entre estudantes (durante a realização do projeto), até a apresentação dos resultados do trabalho \textbf{colaborativo}.
A professora ainda trouxe o exemplo do PCI realizado entre Fatec Guaratinguetá (Brasil) e Uniminuto (Colômbia) em três edições.
No segundo semestre de 2020, o PCI se dedicou a estudar os jargões em diversos \textbf{contextos} comunicativos.
No primeiro semestre de 2021, os grupos mistos de brasileiros e colombianos \textbf{analisaram} mensagens publicitárias e no segundo semestre desse ano, outra \textbf{turma} \textbf{comparou} publicações da empresa Burger King em suas páginas de Facebook no Brasil e na Colômbia. “A metodologia COIL como \textbf{estratégia} intercultural \textbf{permite} o fortalecimento das \textbf{competências} interculturais e do \textbf{aprendizado} de \textbf{idiomas} de \textbf{forma} alternativa e prática”, concluiu Regiane.

% matched lemmas: analisar, aprendizado, colaborativo, comparar, compartilhar, competência, contexto, desenvolvimento, encontro, específico, estratégia, forma, idioma, permitir, que, ser, turma
\end{textsample}
